\documentclass[10pt,a4paper]{article}

\usepackage[a4paper,top=13mm,bottom=15mm,left=15mm,right=15mm]{geometry}
\usepackage{fontspec}
\usepackage[table]{xcolor}
\usepackage{tabularx}
\usepackage{array}
\usepackage{booktabs}
\usepackage{enumitem}
\usepackage[most]{tcolorbox}
\usepackage{tikz}
\usepackage{fancyhdr}
\usepackage{lastpage}
\usepackage{microtype}
\usepackage{ragged2e}

\IfFontExistsTF{Arial}{
  \setmainfont{Arial}
  \setsansfont{Arial}
}{
  \setmainfont{TeX Gyre Heros}
  \setsansfont{TeX Gyre Heros}
}

\definecolor{AIEPRed}{HTML}{D62027}
\definecolor{AIEPRedSoft}{HTML}{FDEBEC}
\definecolor{AIEPNavy}{HTML}{082B4C}
\definecolor{AIEPNavyMid}{HTML}{17476F}
\definecolor{AIEPBlueSoft}{HTML}{EDF4FA}
\definecolor{AIEPInk}{HTML}{1F2933}
\definecolor{AIEPMuted}{HTML}{536171}
\definecolor{AIEPLine}{HTML}{D8E1EA}
\definecolor{AIEPSoft}{HTML}{F6F8FA}
\definecolor{AIEPGold}{HTML}{C59A32}
\definecolor{AIEPGoldSoft}{HTML}{FBF5E6}
\definecolor{AIEPGreen}{HTML}{26734D}
\definecolor{AIEPGreenSoft}{HTML}{EAF5EF}

\pagestyle{fancy}
\fancyhf{}
\renewcommand{\headrulewidth}{0pt}
\fancyfoot[L]{\footnotesize\textcolor{AIEPMuted}{GeoGreen Escolar · AIEP Osorno}}
\fancyfoot[C]{\footnotesize\textcolor{AIEPMuted}{Cronograma operativo 2026}}
\fancyfoot[R]{\footnotesize\textcolor{AIEPMuted}{\thepage\ de \pageref{LastPage}}}

\setlength{\parindent}{0pt}
\setlength{\parskip}{4pt}
\setlength{\tabcolsep}{4.5pt}
\setlist[itemize]{leftmargin=*,topsep=2pt,itemsep=2pt}
\RaggedRight
\hyphenpenalty=8000
\exhyphenpenalty=8000
\sloppy

\newcolumntype{Y}{>{\RaggedRight\arraybackslash}X}
\newcolumntype{L}[1]{>{\RaggedRight\arraybackslash}p{#1}}

\newtcolorbox{hero}{
  enhanced,
  colback=AIEPNavy,
  colframe=AIEPNavy,
  boxrule=0pt,
  arc=2mm,
  left=8mm,
  right=8mm,
  top=5.5mm,
  bottom=4.5mm,
  before skip=0pt,
  after skip=4mm,
  borderline south={1.2mm}{0pt}{AIEPRed}
}

\newtcolorbox{factbox}[2][]{
  enhanced,
  colback=white,
  colframe=AIEPLine,
  boxrule=0.55pt,
  arc=1.6mm,
  left=3.5mm,
  right=3.5mm,
  top=2.8mm,
  bottom=2.5mm,
  before skip=0pt,
  after skip=0pt,
  title={\scriptsize\bfseries\color{AIEPMuted}\MakeUppercase{#2}},
  attach title to upper={\par\vspace{1.2mm}},
  #1
}

\newtcolorbox{phasebox}[3][]{
  enhanced,
  colback=#2,
  colframe=#3,
  boxrule=0.65pt,
  arc=1.6mm,
  left=3.5mm,
  right=3.5mm,
  top=3mm,
  bottom=3mm,
  before skip=0pt,
  after skip=0pt,
  #1
}

\newtcolorbox{chainbox}{
  enhanced,
  colback=AIEPSoft,
  colframe=AIEPRed,
  boxrule=0.7pt,
  arc=1.6mm,
  left=4mm,
  right=4mm,
  top=2.5mm,
  bottom=2.5mm,
  before skip=3mm,
  after skip=0pt
}

\newtcolorbox{sectionbar}[2][]{
  enhanced,
  colback=#2,
  colframe=#2,
  boxrule=0pt,
  arc=1.2mm,
  left=4mm,
  right=4mm,
  top=2.2mm,
  bottom=2.2mm,
  before skip=2.5mm,
  after skip=1mm,
  #1
}

\newcommand{\stepnum}[2]{%
  \begin{tikzpicture}[baseline=(n.base)]
    \node[circle,fill=#1,minimum size=6.5mm,inner sep=0pt,text=white,font=\bfseries\small] (n) {#2};
  \end{tikzpicture}%
}

\newcommand{\ownerred}[1]{{\scriptsize\bfseries\color{AIEPRed}\MakeUppercase{#1}}}
\newcommand{\ownernavy}[1]{{\scriptsize\bfseries\color{AIEPNavyMid}\MakeUppercase{#1}}}
\newcommand{\ownergreen}[1]{{\scriptsize\bfseries\color{AIEPGreen}\MakeUppercase{#1}}}
\newcommand{\datecell}[1]{\textbf{\textcolor{AIEPNavy}{#1}}}
\newcommand{\product}[1]{\space{\footnotesize\textcolor{AIEPMuted}{\textbf{Producto:} #1}}}

\begin{document}
\fontsize{10.2}{12.6}\selectfont

\begin{hero}
{\small\bfseries\color{white!78}\MakeUppercase{Cronograma ejecutivo · 2026}\par}
\vspace{2.5mm}
{\fontsize{27}{31}\selectfont\bfseries\color{white}GeoGreen Escolar\par}
\vspace{1mm}
{\Large\color{white!82}Talleres, mentorías y cierre del programa\par}

\vspace{4mm}
\begin{tabularx}{\linewidth}{@{}Y Y Y@{}}
{\scriptsize\bfseries\color{white!65}LANZAMIENTO\par}
{\normalsize\bfseries\color{white}22 de junio}
&
{\scriptsize\bfseries\color{white!65}INICIO DE TALLERES\par}
{\normalsize\bfseries\color{white}17 de agosto}
&
{\scriptsize\bfseries\color{white!65}EVENTO FINAL\par}
{\normalsize\bfseries\color{white}5 de octubre}
\end{tabularx}
\end{hero}

{\Large\bfseries\color{AIEPNavy}Resumen ejecutivo\par}
\vspace{1.5mm}
GeoGreen Escolar articula sostenibilidad, comprensión de los residuos,
tecnología y comunicación. Acompaña a los equipos desde la identificación de un
problema ambiental hasta la presentación de una propuesta ante jurado.

\vspace{3mm}
\begin{tabularx}{\textwidth}{@{}Y@{\hspace{3mm}}Y@{}}
\begin{factbox}[height=21mm]{Socio comunitario}
\textbf{Instituto Comercial Liceo Bicentenario}\\[-0.3mm]
\textcolor{AIEPMuted}{Osorno}
\end{factbox}
&
\begin{factbox}[height=21mm]{Cobertura}
\textbf{Aproximadamente 60 estudiantes}\\[-0.3mm]
\textcolor{AIEPMuted}{2 bloques · 5 equipos de 6 por bloque}
\end{factbox}
\end{tabularx}

\vspace{2mm}
\begin{tabularx}{\textwidth}{@{}Y@{\hspace{3mm}}Y@{}}
\begin{factbox}[height=21mm]{Trayecto formativo}
\textbf{4 talleres + 4 mentorías}\\[-0.3mm]
\textcolor{AIEPMuted}{Trabajo de los equipos entre sesiones}
\end{factbox}
&
\begin{factbox}[height=21mm]{Cierre}
\textbf{Competencia final y premiación}\\[-0.3mm]
\textcolor{AIEPMuted}{Cierre interno durante la semana del 12 de octubre}
\end{factbox}
\end{tabularx}

\vspace{4mm}
{\Large\bfseries\color{AIEPNavy}Ruta formativa\par}
\vspace{2.5mm}
\begin{tabularx}{\textwidth}{@{}Y@{\hspace{2.5mm}}Y@{\hspace{2.5mm}}Y@{\hspace{2.5mm}}Y@{}}
\begin{phasebox}[height=35mm]{AIEPRedSoft}{AIEPRed}
\stepnum{AIEPRed}{1}\hspace{1.5mm}\textbf{\color{AIEPRed}EXPLORAR}\par
\vspace{2mm}
\textbf{Talleres 1 y 2}\par
\small Problema ambiental, residuos y materiales.
\end{phasebox}
&
\begin{phasebox}[height=35mm]{AIEPBlueSoft}{AIEPNavyMid}
\stepnum{AIEPNavy}{2}\hspace{1.5mm}\textbf{\color{AIEPNavy}PROTOTIPAR}\par
\vspace{2mm}
\textbf{Talleres 3 y 4}\par
\small Sensores, datos y software útil.
\end{phasebox}
&
\begin{phasebox}[height=35mm]{AIEPGoldSoft}{AIEPGold}
\stepnum{AIEPGold}{3}\hspace{1.5mm}\textbf{\color{AIEPNavy}DESARROLLAR}\par
\vspace{2mm}
\textbf{Mentorías 1 a 4}\par
\small Validación, avances y preparación.
\end{phasebox}
&
\begin{phasebox}[height=35mm]{AIEPGreenSoft}{AIEPGreen}
\stepnum{AIEPGreen}{4}\hspace{1.5mm}\textbf{\color{AIEPGreen}PRESENTAR}\par
\vspace{2mm}
\textbf{Evento final}\par
\small Propuesta, jurado y premiación.
\end{phasebox}
\end{tabularx}

\vspace{4mm}
{\Large\bfseries\color{AIEPNavy}Responsables principales\par}
\vspace{2mm}
\renewcommand{\arraystretch}{1.32}
\begin{tabularx}{\textwidth}{@{}L{55mm}Y@{}}
\rowcolor{AIEPRedSoft}
{\small\bfseries\color{AIEPRed}DESARROLLO SOCIAL} &
Talleres 1 y 2 · Mentorías 1 y 4 \\[1.3mm]
\rowcolor{AIEPBlueSoft}
{\small\bfseries\color{AIEPNavyMid}INGENIERÍA, ENERGÍA Y TECNOLOGÍA (I/E/T)} &
Talleres 3 y 4 · Mentorías 2 y 3 \\[1.3mm]
\rowcolor{AIEPGreenSoft}
{\small\bfseries\color{AIEPGreen}EQUIPO CONJUNTO} &
Hito comunicacional · Evento final · Cierre interno \\
\end{tabularx}
\renewcommand{\arraystretch}{1}

\begin{chainbox}
\textbf{\color{AIEPRed}CADA PRODUCTO ALIMENTA LA ETAPA SIGUIENTE}\par
\small Los equipos desarrollan entre sesiones; las mentorías revisan, orientan
y destraban avances sin reemplazar su trabajo.
\end{chainbox}

\newpage

{\small\bfseries\color{AIEPRed}\MakeUppercase{GeoGreen Escolar · Programación 2026}\par}
\vspace{1.5mm}
{\fontsize{23}{27}\selectfont\bfseries\color{AIEPNavy}Calendario detallado\par}
\vspace{1mm}
\textcolor{AIEPMuted}{Secuencia cronológica de actividades, propósito y producto esperado.}
\vspace{2mm}

\begin{sectionbar}{AIEPNavy}
{\bfseries\color{white}1 · PREPARACIÓN Y LANZAMIENTO}
\end{sectionbar}
\renewcommand{\arraystretch}{1.18}
\begin{tabularx}{\textwidth}{@{}L{25mm}L{39mm}Y@{}}
\rowcolor{AIEPNavy}
\textbf{\color{white}Fecha} & \textbf{\color{white}Actividad} & \textbf{\color{white}Propósito y producto} \\
\rowcolor{AIEPSoft}
\datecell{Lun. 22 jun.} & \textbf{Lanzamiento} &
Presentar GeoGreen Escolar y abrir la coordinación operativa con el establecimiento.
\product{Acuerdos iniciales.} \\[0.8mm]
\rowcolor{white}
\datecell{Jun.–jul.} & \textbf{Preparación interna} &
Revisar materiales, sensores, fichas, presentaciones y condiciones de ejecución.
\product{Lista de materiales.} \\
\end{tabularx}

\begin{sectionbar}{AIEPRed}
{\bfseries\color{white}2 · TALLERES · CONSTRUCCIÓN DE LA PROPUESTA}
\end{sectionbar}
\begin{tabularx}{\textwidth}{@{}L{25mm}L{39mm}Y@{}}
\rowcolor{AIEPRedSoft}
\datecell{Lun. 17 ago.} &
\textbf{Taller 1}\par\ownerred{Desarrollo Social} &
Observar residuos, hábitos y situaciones concretas del entorno escolar.
\product{Problema ambiental definido.} \\[0.7mm]
\rowcolor{white}
\datecell{Mar. 18 ago.} &
\textbf{Taller 2}\par\ownerred{Desarrollo Social} &
Comprender materiales, reciclabilidad, separación y contaminación.
\product{Ficha de análisis del residuo o material.} \\[0.7mm]
\rowcolor{AIEPBlueSoft}
\datecell{Lun. 24 ago.} &
\textbf{Taller 3}\par\ownernavy{I/E/T} &
Relacionar el problema con sensores, Arduino, visualización y alertas.
\product{Idea tecnológica inicial.} \\[0.7mm]
\rowcolor{white}
\datecell{Mar. 25 ago.} &
\textbf{Taller 4}\par\ownernavy{I/E/T} &
Convertir datos simulados en estados, interfaces, alertas y decisiones.
\product{Prototipo digital provisional y evidencia de prueba.} \\
\end{tabularx}

\begin{sectionbar}{AIEPGold}
{\bfseries\color{white}3 · MENTORÍAS · REVISIÓN Y DESARROLLO}
\end{sectionbar}
\begin{tabularx}{\textwidth}{@{}L{25mm}L{39mm}Y@{}}
\rowcolor{AIEPGoldSoft}
\datecell{Lun. 31 ago.} &
\textbf{Mentoría 1}\par\ownerred{Desarrollo Social} &
Revisar problema, usuario, contexto y primera solución.
\product{Problema validado.} \\[0.5mm]
\rowcolor{white}
\datecell{Lun. 7 sep.} &
\textbf{Mentoría 2}\par\ownernavy{I/E/T} &
Ajustar solución, recursos, materiales y roles de trabajo.
\product{Solución y recursos definidos.} \\[0.5mm]
\rowcolor{AIEPRedSoft}
\datecell{Sem. 14 sep.} &
\textbf{Sin actividades} &
Semana de Fiestas Patrias.
\product{Sin sesión programada.} \\[0.5mm]
\rowcolor{white}
\datecell{Lun. 21 sep.} &
\textbf{Mentoría 3}\par\ownernavy{I/E/T} &
Revisar maqueta, simulación, diagrama, prototipo o evidencia de avance.
\product{Evidencia revisada.} \\[0.5mm]
\rowcolor{AIEPGoldSoft}
\datecell{Lun. 28 sep.} &
\textbf{Mentoría 4 y ensayo}\par\ownerred{Desarrollo Social} &
Afinar guion, soporte visual, roles y tiempos de presentación.
\product{Presentación preparada.} \\
\end{tabularx}

\begin{sectionbar}{AIEPGreen}
{\bfseries\color{white}4 · COMUNICACIÓN, EVENTO FINAL Y CIERRE}
\end{sectionbar}
\begin{tabularx}{\textwidth}{@{}L{25mm}L{39mm}Y@{}}
\rowcolor{AIEPGreenSoft}
\datecell{Vie. 2 oct.} & \textbf{Hito comunicacional} &
Comunicar el cierre próximo del proceso.
\product{Difusión breve.} \\[0.5mm]
\rowcolor{white}
\datecell{Lun. 5 oct.} & \textbf{Competencia final} &
Presentar ante jurado, recibir retroalimentación y realizar la premiación.
\product{Propuestas presentadas y registro.} \\[0.5mm]
\rowcolor{AIEPGreenSoft}
\datecell{Sem. 12 oct.} & \textbf{Cierre interno} &
Consolidar evidencias, aprendizajes y recomendaciones del proceso.
\product{Informe y materiales reutilizables.} \\
\end{tabularx}

\vfill
{\footnotesize\textcolor{AIEPMuted}{%
\textbf{Lectura del calendario:} cada producto constituye la entrada de la etapa
siguiente y permite seguir la progresión completa del proyecto.}}

\end{document}
